\documentclass{article}
\usepackage[UTF8]{ctex}
\usepackage{amsmath,amssymb,amsfonts}
\usepackage{enumitem}
\usepackage{geometry}
\geometry{a4paper, margin=1in}
% === 错题本排版优化（不参与题目标识） ===
\setlength{\parindent}{0pt}
\setlength{\parskip}{0.35em plus 0.10em minus 0.08em}
\setlength{\emergencystretch}{3em}
\setlist{itemsep=0.25em, topsep=0.40em, parsep=0pt}
\allowdisplaybreaks[2]
% === 错题本排版优化结束 ===
\begin{document}
\section*{太原理工大学2026年数学分析真题}
\section*{68．太原理工大学2026年数学分析试卷}
\begin{itemize}

\subsection{题目1}
  \item[1．] 求极限
\[
\lim _{n \rightarrow \infty}\left(\frac{\sin \frac{\pi}{n}}{n+2026}+\frac{\sin \frac{2 \pi}{n}}{n+\frac{2026}{2}}+\cdots+\frac{\sin \frac{n \pi}{n}}{n+\frac{2026}{n}}\right) .
\]

\subsection{题目2}
  \item[2．] 求极限 \(\lim _{x \rightarrow 0} \frac{\tan (\tan x)-\sin (\sin x)}{x^{2} \arctan 2 x}\) ．
\end{itemize}

\subsection{题目3}
3．计算不定积分 \(\int \frac{\cos ^{2} x}{\left(a^{2} \sin ^{2} x+b^{2} \cos ^{2} x\right)^{2}} \mathrm{~d} x\) ．

\subsection{题目4}
4．计算曲线 \(\left(\frac{x^{2}}{a^{2}}+\frac{y^{2}}{b^{2}}\right)^{2}=x^{2}+y^{2}(a, b>0)\) 围成的平面图形的面积．

\subsection{题目5}
5．计算曲线积分

\[
\oint_{L}(z-y) \mathrm{d} x+(x-z) \mathrm{d} y+(x-y) \mathrm{d} z
\]

其中 \(L\) 是 \(x^{2}+y^{2}=4\) 与 \(x-y+z=1\) 的交线，从 \(z\) 轴正向往 \(z\) 轴负向看是顺时针方向。

\subsection{题目6}
6．计算第二型曲面积分

\[
\iint_{S} 2\left(1-x^{2}\right) \mathrm{d} y \mathrm{~d} z+8 x y \mathrm{~d} z \mathrm{~d} x-4 x z \mathrm{~d} x \mathrm{~d} y
\]

其中 \(S\) 是 \(x=5^{y}(0 \leq y \leq 1)\) 绕 \(x\) 轴形成的曲面的外侧。

\subsection{题目7}
7．计算 \(\int_{0}^{+\infty} \frac{\arctan x}{x\left(1+x^{2}\right)} \mathrm{d} x\) ．

\subsection{题目8}
8．判断反常积分 \(\int_{0}^{+\infty} \frac{x \ln x}{x^{p}+1}(p>0)\) 的敛散性．

\subsection{题目9}
9．设 \(f(x)\) 是 \(2 \pi\) 为周期的函数，\(f(x)=x(\pi-x), x \in[0, \pi]\) ．

\begin{itemize}
  \item[（1）] 求 \(f(x)\) 的余弦级数和正弦级数展开式．
  \item[（2）] 求 \(\sum_{n=1}^{\infty} \frac{1}{n^{2}}\) 和 \(\sum_{n=1}^{\infty} \frac{(-1)^{n-1}}{(2 n-1)^{3}}\) ．
\end{itemize}

\subsection{题目10}
10．设函数 \(f(x)\) 在 \([-1,1]\) 上连续，令 \(\varphi_{n}(x)=\left\{\begin{array}{ll}(1-x)^{n}, & 0 \leq x \leq 1 ; \\ (1+x)^{n}, & -1 \leq x \leq 0 .\end{array}\right.\) 证明：

\[
\lim _{n \rightarrow \infty} \frac{n+1}{2} \int_{-1}^{1} f(x) \varphi_{n}(x) \mathrm{d} x=f(0)
\]

\subsection{题目11}
11．数列 \(\left\{a_{n}\right\}\) 满足 \(a_{0}=1, a_{1}=1, a_{n}=a_{n-1}+2 a_{n-2}(n \geq 2)\) ．

\begin{itemize}
  \item[（1）] 证明：\(\left(\frac{5}{3}\right)^{n-1} \leq a_{n} \leq 3^{n-1}\) ．
  \item[（2）] 证明：\(\sum_{n=1}^{\infty} \frac{1}{\ln a_{n}}\) 发散，\(\sum_{n=1}^{\infty} \frac{1}{a_{n}}\) 收敛。
  \item[（3）] 若 \(\lim _{n \rightarrow \infty} \frac{a_{n}}{a_{n+1}}\) 存在，求幂级数 \(\sum_{n=1}^{\infty} a_{n} x^{n}\) 的收敛半径以及收敛区间内的和函数．
\end{itemize}

\subsection{题目12}
12．解答如下问题：

\begin{itemize}
  \item[（1）] 设 \(f(x), g(x)\) 都是凸函数，\(f(x) g(x)\) 和 \(f(g(x))\) 是否是凸函数？若不是，给出反例．对于 \(f(g(x))\) ，如何加一个充分条件，使之成为凸函数？给出证明。
  \item[（2）] 证明：实数域上的有界凸函数必是常值函数．
\end{itemize}
\end{document}
